\documentclass[a4paper,12pt]{article}

\usepackage[T2A]{fontenc}
\usepackage[utf8]{inputenc}
\usepackage[russian]{babel}
\usepackage{amsmath,amssymb,mathtools}
\usepackage{helvet}
\renewcommand{\familydefault}{\sfdefault}
\usepackage{icomma}
\usepackage[a4paper,margin=20mm,headheight=15pt]{geometry}
\usepackage{microtype}
\usepackage{tikz}
\usetikzlibrary{calc,arrows.meta,positioning,shapes.geometric}
\usepackage{tabularx,booktabs}
\usepackage{enumitem}
\usepackage{hyperref}
\usepackage{parskip}
\usepackage{fancyhdr}
\usepackage{setspace}
\usepackage{xcolor}

\definecolor{accent}{HTML}{008080}
\definecolor{darkaccent}{HTML}{004D4D}
\definecolor{textgray}{HTML}{333333}
\definecolor{softgray}{HTML}{666666}

\hypersetup{
  colorlinks=true,
  linkcolor=darkaccent,
  urlcolor=darkaccent
}

\color{textgray}
\onehalfspacing
\setlength{\parindent}{0pt}
\setlength{\emergencystretch}{2em}

\setlist[itemize]{
  leftmargin=1.5em,
  itemsep=0.2em,
  topsep=0.35em
}

\setlist[enumerate]{
  leftmargin=1.8em,
  itemsep=0.35em,
  topsep=0.35em
}

\pagestyle{fancy}
\fancyhf{}
\lhead{\small Движение с постоянным ускорением}
\rhead{\small Володя Хачатурян}
\cfoot{\small \thepage}
\renewcommand{\headrulewidth}{0.3pt}

\newcommand{\sectionline}{%
  \vspace{0.1em}
  \noindent\textcolor{accent}{\rule{\linewidth}{0.5pt}}
  \vspace{0.45em}
}

\newcommand{\key}[1]{\textcolor{darkaccent}{\textbf{#1}}}
\newcommand{\ms}{\,\text{м/с}}
\newcommand{\mss}{\,\text{м/с}^{2}}
\newcommand{\metre}{\,\text{м}}
\newcommand{\second}{\,\text{с}}

\tikzset{
  flowstart/.style={
    ellipse,
    draw=darkaccent,
    line width=0.9pt,
    minimum width=34mm,
    minimum height=10mm,
    align=center,
    font=\small
  },
  flowaction/.style={
    rounded corners=3pt,
    draw=darkaccent,
    line width=0.9pt,
    text width=63mm,
    minimum height=13mm,
    align=center,
    font=\small,
    inner sep=4pt
  },
  flowdecision/.style={
    diamond,
    aspect=2.15,
    draw=darkaccent,
    line width=0.9pt,
    text width=39mm,
    align=center,
    font=\small,
    inner sep=2pt
  },
  flowarrow/.style={
    -{Stealth[length=3mm,width=2mm]},
    line width=0.9pt,
    draw=darkaccent
  }
}

\begin{document}

% =========================================================
% ТИТУЛЬНЫЙ ЛИСТ
% =========================================================

\begin{titlepage}
\thispagestyle{empty}
\centering

\vspace*{24mm}

{\Large\textcolor{darkaccent}{\textbf{Чек-лист}}\par}

\vspace{7mm}

{\huge\textbf{Равноускоренное прямолинейное движение}\par}

\vspace{3mm}

{\Large\textbf{Координата, скорость, перемещение и путь за отдельную секунду}\par}

\vspace{18mm}

{\large Ученик: \textbf{Володя Хачатурян}\par}
\vspace{3mm}
{\large Подготовка: \textbf{ОГЭ по физике}\par}
\vspace{3mm}
{\large Тренировочный блок: \textbf{Путь Б}\par}

\vfill

{\large 23 сентября 2026 г.\par}

\vspace{6mm}

{\large\textbf{Лёвин Артём Александрович эксклюзивно для Володи Хачатуряна}\par}

\vspace{16mm}

\begin{tikzpicture}[remember picture,overlay]
  \draw[accent,line width=1.1pt]
  ([xshift=18mm,yshift=17mm]current page.south west)
  .. controls
  ([xshift=70mm,yshift=28mm]current page.south west)
  and
  ([xshift=-70mm,yshift=8mm]current page.south east)
  ..
  ([xshift=-18mm,yshift=18mm]current page.south east);
\end{tikzpicture}

\end{titlepage}

% =========================================================
% ОСНОВНОЙ ЧЕК-ЛИСТ
% =========================================================

\section*{1. Три величины, которые нельзя путать}
\sectionline

\begin{itemize}
  \item \key{Координата $x$} показывает положение тела на оси $Ox$.

  \item \key{Перемещение $\Delta x$} за промежуток времени:
  \[
    \boxed{\Delta x=x_{\text{кон}}-x_{\text{нач}}}.
  \]
  Перемещение может быть положительным, отрицательным или равным нулю.

  \item \key{Путь $s$} --- длина пройденного участка траектории, поэтому
  \[
    s\ge 0.
  \]
\end{itemize}

Если тело на рассматриваемом промежутке не меняет направление движения, то
\[
  s=|\Delta x|.
\]
Если направление меняется, путь нужно находить по отдельным участкам движения.

\section*{2. Базовые формулы движения с постоянным ускорением}
\sectionline

Для движения вдоль оси $Ox$ при постоянной проекции ускорения $a_x$:

\[
\boxed{x=x_0+v_{0x}t+\frac{a_xt^2}{2}}
\]

\[
\boxed{v_x=v_{0x}+a_xt}
\]

\[
\boxed{\Delta x=v_{0x}t+\frac{a_xt^2}{2}}
\]

где $x_0$ --- начальная координата, $v_{0x}$ --- начальная проекция скорости,
$a_x$ --- проекция ускорения, $t$ --- время.

\textbf{Главное правило знаков.}
Знак входит в значение коэффициента. Например, если
\[
  x=6-3t,
\]
то
\[
  v_{0x}=-3\ms,
  \qquad
  a_x=0.
\]

\section*{3. Как читать уравнение координаты}
\sectionline

Если закон движения записан как
\[
  x=A+Bt+Ct^2,
\]
то сравниваем его с
\[
  x=x_0+v_{0x}t+\frac{a_x}{2}t^2.
\]
Отсюда
\[
\boxed{x_0=A,\qquad v_{0x}=B,\qquad a_x=2C.}
\]

Особенно важно помнить: коэффициент при $t^2$ равен $a_x/2$, поэтому ускорение в два раза больше этого коэффициента.

\subsection*{Пример 1. В уравнении отсутствуют слагаемые}

Пусть
\[
  x=6t^2.
\]
Полная запись имеет вид
\[
  x=0+0\cdot t+6t^2.
\]
Следовательно,
\[
  x_0=0,
  \qquad
  v_{0x}=0,
  \qquad
  \frac{a_x}{2}=6,
\]
поэтому
\[
  \boxed{a_x=12\mss.}
\]

А для уравнения
\[
  x=6-3t
\]
имеем
\[
  \boxed{x_0=6\metre,\qquad v_{0x}=-3\ms,\qquad a_x=0.}
\]

\subsection*{Пример 2. Координата в заданный момент}

Пусть
\[
  x=8t-2t^2.
\]
Через $3\second$:
\[
  x(3)=8\cdot3-2\cdot3^2=24-18=6\metre.
\]

\section*{4. Как найти перемещение за несколько секунд}
\sectionline

Если известен закон координаты, надёжнее всего использовать
\[
\boxed{\Delta x=x(t_2)-x(t_1).}
\]

Например, при
\[
  x=6-4t+t^2
\]
за первые $5\second$:
\[
  x(0)=6\metre,
\]
\[
  x(5)=6-4\cdot5+5^2=11\metre,
\]
поэтому
\[
  \boxed{\Delta x=11-6=5\metre.}
\]

\textbf{Самопроверка.}
Если требуется перемещение «за первые $T$ секунд», нужно сравнить координаты в моменты $0$ и $T$.

\section*{5. Перемещение за $n$-ю секунду}
\sectionline

Фраза «за $n$-ю секунду» означает промежуток
\[
  [n-1;\,n]\ \text{с}.
\]
Например, четвёртая секунда --- это промежуток от $3\second$ до $4\second$.

Пусть $\tau=1\second$ --- длительность одной секунды. Тогда
\[
  \Delta x_n=x(n\tau)-x((n-1)\tau).
\]
Подставим закон движения:
\begin{align*}
\Delta x_n
&=
\left(
  x_0+v_{0x}n\tau+\frac{a_xn^2\tau^2}{2}
\right)
-
\left(
  x_0+v_{0x}(n-1)\tau+
  \frac{a_x(n-1)^2\tau^2}{2}
\right)
\\[1mm]
&=v_{0x}\tau+
\frac{a_x\tau^2}{2}
\left(n^2-(n-1)^2\right)
\\[1mm]
&=v_{0x}\tau+
\frac{a_x\tau^2}{2}(2n-1).
\end{align*}

Итак,
\[
\boxed{
\Delta x_n=
 v_{0x}\tau+
 \frac{a_x\tau^2}{2}(2n-1),
\qquad \tau=1\second.
}
\]

Для вычислений можно сразу подставить $\tau=1\second$:
\[
\boxed{
\Delta x_n=
 v_{0x}\cdot 1\second+
 \frac{a_x(1\second)^2}{2}(2n-1).
}
\]
Так сохраняются и физический смысл, и размерность формулы.

\subsection*{Пример 3. Четвёртая и седьмая секунды}

Начальная скорость тела равна
\[
  v_{0x}=2\ms,
\]
а путь за четвёртую секунду составляет
\[
  s_4=2{,}7\metre.
\]
В рассматриваемой задаче движение идёт в положительном направлении, поэтому на этом интервале
\[
  s_4=\Delta x_4.
\]

По формуле для четвёртой секунды:
\[
  2{,}7=2+\frac{a_x}{2}(2\cdot4-1)
  =2+\frac{7a_x}{2}.
\]
Следовательно,
\[
  \frac{7a_x}{2}=0{,}7,
  \qquad
  a_x=0{,}2\mss.
\]

Теперь
\[
  \Delta x_7=
  2+\frac{0{,}2}{2}(2\cdot7-1)
  =2+0{,}1\cdot13
  =3{,}3\metre.
\]
Так как скорость остаётся положительной, путь за седьмую секунду также равен
\[
  \boxed{s_7=3{,}3\metre.}
\]

\section*{6. Когда путь и перемещение за секунду различаются}
\sectionline

Чтобы понять, можно ли использовать $s=|\Delta x|$, проверьте знак скорости на рассматриваемом интервале.

Если
\[
  v_x=v_{0x}+a_xt
\]
обращается в ноль внутри интервала $[t_1;t_2]$, тело меняет направление. Тогда нужно найти момент остановки $t_*$ из уравнения
\[
  v_x(t_*)=0
\]
и сложить пути до и после остановки:
\[
\boxed{
  s=|x(t_*)-x(t_1)|+|x(t_2)-x(t_*)|.
}
\]

\section*{7. Как читать уравнение скорости}
\sectionline

Если
\[
  v_x=D+Et,
\]
то сравнение с
\[
  v_x=v_{0x}+a_xt
\]
даёт
\[
\boxed{v_{0x}=D,\qquad a_x=E.}
\]

\subsection*{Пример 4. Когда скорость становится равной нулю}

Пусть
\[
  x=12t-2t^2.
\]
Из закона координаты:
\[
  v_{0x}=12\ms,
  \qquad
  \frac{a_x}{2}=-2,
  \qquad
  a_x=-4\mss.
\]
Поэтому
\[
  v_x=12-4t.
\]
В момент остановки
\[
  v_x=0,
\]
значит
\[
  12-4t=0,
  \qquad
  \boxed{t=3\second.}
\]

\subsection*{Пример 5. Из скорости получить перемещение}

Пусть
\[
  v_x=3-2t.
\]
Тогда
\[
  v_{0x}=3\ms,
  \qquad
  a_x=-2\mss.
\]
Следовательно,
\[
  \Delta x=v_{0x}t+\frac{a_xt^2}{2}
  =3t-t^2.
\]
Итак,
\[
  \boxed{\Delta x=3t-t^2.}
\]

\section*{8. Перевод формулировок задачи на язык физики}
\sectionline

\renewcommand{\arraystretch}{1.3}
\begin{table}[ht]
\centering
\begin{tabularx}{\linewidth}{@{}>{\raggedright\arraybackslash}X>{\raggedright\arraybackslash}X@{}}
\toprule
\textbf{Фраза в условии} & \textbf{Что нужно увидеть} \\
\midrule
«Зависимость координаты от времени описывается уравнением...»
& Дано уравнение $x(t)$. \\
«Движение описывается уравнением $x=\ldots$»
& Дано уравнение координаты $x(t)$. \\
«В какой момент времени...»
& Искомая величина --- $t$. \\
«Проекция скорости на ось $Ox$ равна нулю»
& Использовать условие $v_x=0$. \\
«Перемещение за первые $5$ секунд»
& Найти $x(5)-x(0)$. \\
«За четвёртую секунду»
& Рассмотреть интервал от $3$ до $4$ с. \\
\bottomrule
\end{tabularx}
\end{table}

\section*{9. Типичные ошибки}
\sectionline

\begin{enumerate}
  \item Путать величину \textbf{за четыре секунды} с величиной \textbf{за четвёртую секунду}.

  \item Терять знак коэффициента. В выражении $x=6-3t$ начальная скорость равна $-3\ms$.

  \item Принимать коэффициент при $t^2$ непосредственно за ускорение. В законе координаты этот коэффициент равен $a_x/2$.

  \item Добавлять $x_0$ в формулу перемещения. Начальная координата входит в формулу координаты, но сокращается при вычислении $\Delta x$.

  \item Считать путь равным перемещению при смене направления движения.

  \item Подставлять числа до того, как определена искомая физическая величина и выбрана формула.

  \item Записывать число без его знака или без единицы измерения.
\end{enumerate}

% =========================================================
% БЛОК-СХЕМА
% =========================================================

\clearpage
\thispagestyle{empty}

\begin{center}
{\Large\textcolor{darkaccent}{\textbf{Алгоритм: как разобрать уравнение движения}}\par}

\vspace{4mm}

\begin{tikzpicture}[node distance=6mm and 16mm]

\node[flowstart] (start) {Начало};

\node[flowdecision,below=of start] (isx)
{Дано уравнение координаты $x(t)$?};

\node[flowaction,below=11mm of isx,xshift=-41mm] (xform)
{Сравните с $x=x_0+v_{0x}t+\dfrac{a_xt^2}{2}$.\\
Найдите $x_0$, $v_{0x}$, $a_x$.};

\node[flowaction,below=11mm of isx,xshift=41mm] (vform)
{Тогда дано $v_x(t)$. Сравните с $v_x=v_{0x}+a_xt$.\\
Найдите $v_{0x}$, $a_x$.};

\node[flowaction,below=42mm of isx] (solve)
{Переведите вопрос в обозначения, выберите формулу и выполните вычисления, сохраняя знаки.};

\node[flowdecision,below=of solve] (path)
{Искомая величина --- путь $s$?};

\node[flowaction,below=11mm of path,xshift=-43mm] (pathcheck)
{Проверьте знак $v_x$. При смене направления найдите момент остановки и сложите пути по участкам.};

\node[flowaction,below=11mm of path,xshift=43mm] (directcheck)
{Проверьте полученный результат: единицы измерения, знак и физический смысл.};

\node[flowaction,below=40mm of path] (finalcheck)
{Для пути убедитесь, что итог неотрицателен; для проекций сохраните физически значимый знак.};

\node[flowstart,below=of finalcheck] (finish) {Ответ};

\draw[flowarrow] (start) -- (isx);
\draw[flowarrow] (isx) -| node[pos=0.22,above left,font=\small]{да} (xform);
\draw[flowarrow] (isx) -| node[pos=0.22,above right,font=\small]{нет} (vform);
\draw[flowarrow] (xform.south) |- (solve.west);
\draw[flowarrow] (vform.south) |- (solve.east);
\draw[flowarrow] (solve) -- (path);
\draw[flowarrow] (path) -| node[pos=0.22,above left,font=\small]{да} (pathcheck);
\draw[flowarrow] (path) -| node[pos=0.22,above right,font=\small]{нет} (directcheck);
\draw[flowarrow] (pathcheck.south) |- (finalcheck.west);
\draw[flowarrow] (directcheck.south) |- (finalcheck.east);
\draw[flowarrow] (finalcheck) -- (finish);

\end{tikzpicture}
\end{center}

\clearpage

% =========================================================
% ТРЕНИРОВОЧНЫЙ БЛОК
% =========================================================

\section*{Тренировочный блок --- Путь Б}
\sectionline

На каждый день выполните по три задания. Формулу сначала записывайте в общем виде, затем подставляйте численные значения. В ответах обязательно указывайте единицы измерения.

\subsection*{24.09.2026}

\begin{enumerate}[label=\textbf{\arabic*.}]

  \item
  Движение тела описывается уравнением
  \[
    x=5+4t+t^2.
  \]
  Определите $x_0$, $v_{0x}$, $a_x$ и координату тела через $3\second$.

  \item
  Проекция скорости тела задаётся уравнением
  \[
    v_x=7-2t.
  \]
  Определите начальную скорость, ускорение и момент времени, когда $v_x=0$.

  \item
  Тело движется с постоянным ускорением вдоль оси $Ox$. Начальная скорость равна
  \[
    1{,}5\ms,
  \]
  а перемещение за третью секунду равно
  \[
    2{,}5\metre.
  \]
  Найдите ускорение и перемещение за шестую секунду.

\end{enumerate}

\clearpage

\subsection*{25.09.2026}

\begin{enumerate}[label=\textbf{\arabic*.},resume]

  \item
  Движение тела описывается уравнением
  \[
    x=12-3t+2t^2.
  \]
  Определите $x_0$, $v_{0x}$, $a_x$, скорость через $2\second$ и перемещение за первые $2\second$.

  \item
  Проекция скорости тела задаётся уравнением
  \[
    v_x=6-3t.
  \]
  Запишите зависимость перемещения от времени $\Delta x(t)$. Затем найдите момент остановки и перемещение от $t=0$ до этого момента.

  \item
  Движение тела описывается уравнением
  \[
    x=-4+10t-2t^2.
  \]
  Найдите:
  \begin{itemize}
    \item момент, когда скорость равна нулю;
    \item координату тела в этот момент;
    \item перемещение за третью секунду;
    \item путь, пройденный за третью секунду.
  \end{itemize}

\end{enumerate}

% =========================================================
% ОТВЕТЫ
% =========================================================

\clearpage

\section*{Ответы к тренировочному блоку}
\sectionline

\renewcommand{\arraystretch}{1.4}
\begin{table}[ht]
\centering
\begin{tabularx}{\linewidth}{@{}c>{\raggedright\arraybackslash}X@{}}
\toprule
\textbf{№} & \textbf{Ответ} \\
\midrule
1 & $x_0=5\metre$; $v_{0x}=4\ms$; $a_x=2\mss$; $x(3)=26\metre$. \\
2 & $v_{0x}=7\ms$; $a_x=-2\mss$; $t=3{,}5\second$. \\
3 & $a_x=0{,}4\mss$; $\Delta x_6=3{,}7\metre$. \\
4 & $x_0=12\metre$; $v_{0x}=-3\ms$; $a_x=4\mss$; $v_x(2)=5\ms$; $\Delta x=2\metre$. \\
5 & $\Delta x=6t-1{,}5t^2$; $t_{\text{ост}}=2\second$; $\Delta x(2)=6\metre$. \\
6 & $t_{\text{ост}}=2{,}5\second$; $x(2{,}5)=8{,}5\metre$; $\Delta x_3=0$; $s_3=1\metre$. \\
\bottomrule
\end{tabularx}
\end{table}

\end{document}